\chapter{ws-Locally Closed Sets in Topological Spaces}
\graphicspath{{Chapter7/Chapter7Figs/EPS/}{Chapter7/Chapter7Figs/}}

\section{Introduction}\label{sec7.1}

LC-continuity and LC-irresoluteness

Bourbaki \cite{Bourbaki} defined locally closed set in $(\TSP,\tau)$ as intersection of closed and open sets. Stone \cite{Stone} used a new term as FG for locally closed. Further LC-continuity and LC-irresoluteness developed by Ganster and Reilly \cite{Ganster}. Sundaram \cite{Sundaram} developed the knowledge of generalized locally closed sets, the class of GLC-continuous functions, GLC-irresolute functions and analysed some of their topological properties. Arockiarani \cite{Arockiarani} popularized regular generalized locally closed sets and obtained six new classes of generalized continuities using the theory of regular generalized closed sets. Recently Sheik John \cite{Sheik1} introduced three new classes of sets denoted by w-LC$(\TSP, \tau)$, w-LC{*}$(\TSP, \tau)$ and w-LC{*}{*}$(\TSP, \tau)$ each of which contains LC$(\TSP, \tau)$. Also various authors like Park \cite{Park}, Balachandran \cite{Balachandran} have put stremendous effort to the growth of locally closed sets and locally continuous maps.

In the second section of this chapter, we describe three weaker forms of locally closed sets called ws-lc set, ws-lc* set and ws-lc{*}{*} set each of which is weaker than locally closed set and examine some properties of ws-lc set, ws-lc* set and ws-lc{*}{*} set. 

\newpage

In third section of this chapter, we introduce wsLC-continuous, wsLC*-continuous and wsLC{*}{*}-continuous functions Also we introduce ws-submaximal spaces and obtain some of their properties.

In fourth section of this chapter, we define wsLC-irresolute, wsLC*-irresolute and wsLC**-irresolute functions and analyse some of their related properties.

\section{ws-Locally Closed Sets}\label{sec7.2}

\begin{dfn}\label{defi7.2.1}
A subset $\clrD$ of a topological space $(\TSP, \tau)$ is called a weakly semi locally closed set (briefly wslc-set) if $\clrD = \TSS \cap \TSF$ where $\TSS$ is ws-open and $\TSF$ is ws-closed.
\end{dfn}

The set of all weakly semi locally closed set in $(\TSP,\tau)$ is denoted by wsLC$(\TSP, \tau)$.

\begin{exm}\label{exam7.2.1}
Take up $\TSP = \{1,2,3,4\}$ and $\tau = \{\phi,\TSP, \{1\}, \{2\}, \{1,2\},\{1,2,3\}\}$ be a topology on $\TSP$. Then $\ws\TSC(\TSP) = \{\phi, \TSP,\{1\}, \{2\}, \{3\}, \{4\}, \{1,3\}, \{1,4\} \{3,4\}, \{2,3\},\break \{2,4\}, \{1,2,4\},\{1,3,4\},\{2,3,4\}\}$, $\ws\TSO(\TSP) = \{\phi, \TSP, \{1\}, \{2\}, \{3\},\{1,2\}, \{3,4\}, \{1,4\},\break \{2,3\}, \{2,4\}, \{1,2,3\}\{1,2,4\},\{1,3,4\},\{2,3,4\}\}$ and $\ws\LC(\TSP)=\{\phi,\TSP,\{1\},\{2\},\{3\},\break \{4\},\{1,2\},\{1,3\},\{1,4\},\{2,3\},\{2,4\},\{3,4\},\{1,2,3\},\{1,2,4\},\{1,3,4\},\{2,3,4\}\}$
\end{exm}

\begin{rem}\label{rem7.2.1}
\begin{enumerate}[(1)]
\item A subset $\clrD$ of a topological space $(\TSP, \tau)$ is a ws-lc set if its complement is the union of ws-open set and ws-closed set.
\item Each ws-open set of $(\TSP, \tau)$ s a wslc set
\end{enumerate}
\end{rem}

\begin{dfn}\label{defi7.2.2}
A subset $\clrD$ of a space $\TSP$ is termed as $\ws\LC$*-set if $\exists$ ws-open set $\TSS$ and a closed set $\TSF$ of $\TSP$, $\ni: \clrD = \TSS \cap \TSF$.
\end{dfn}

\begin{dfn}\label{defi7.2.3}
A subset $\clrD$ of a space $\TSP$ is termed as $\ws\LC${*}{*}-set if $\exists$ an open set $\TSS$ and ws-closed set $\TSF$ of $\TSP \ni: \clrD = \TSS\cap \TSF$.
\end{dfn}

\begin{thm}\label{thm7.2.1}
If a subset $\clrD$ of $(\TSP, \tau)$ is locally closed then it is $\ws\LC(\TSP, \tau)$, $\ws\LC$*$(\TSP, \tau)$ and $\ws\LC${*}{*}$(\TSP, \tau)$ but inverse is untrue.
\end{thm}

\begin{proof}
Take up $\clrD = \TSU \cap \TSV$, $\forall \TSU$ is open and $\TSV$ is closed in \ref{thm7.2.1} $(\TSP, \tau)$. Since each open set is ws-open and each closed set is ws-closed. By Definition \ref{defi7.2.3}. A subset $\clrD$ of a topological space $(\TSP, \tau)$ is said to be $\ws\LC${*}{*}-set if $\exists$ an open set $\TSS$ and a ws-closed set $\TSF$ of $(\TSP, \tau) \ni: \clrD = \TSS\cap \TSF$., Therefore $\clrD$ is $\ws\LC (\TSP, \tau)$, $\ws\LC$*$(\TSP, \tau)$ and $\ws\LC${*}{*}$(\TSP, \tau)$.

We use example \ref{exam7.2.2} to prove the inverse of theorem is untrue.
\end{proof}

\begin{exm}\label{exam7.2.2}
$\TSP = \{1, 2, 3, 4\}$ and $\tau = \{\phi, \TSP, \{1\}, \{2\}, \{1, 2\}, \{1,2, 3\}\}$ be a topology on $\TSP$. Then $\ws\TSC(\TSP) = \{\phi, \TSP, \{1\}, \{2\}, \{3\}, \{4\}, \{1,3\}, \{1,4\},\{3,4\},\{2,3\}, \{2,4\},\break \{1,2, 4\},\{1,3,4\},\{2,3,4\}\}$, $\ws\TSO(\TSP) = \{\phi, \TSP, \{1\}, \{2\}, \{3\},\{1,2\},\{1,3\} \{3,4\}, \{1,4\},\break \{2,3\}, \{2,4\}, \{1,2,3\},\{1,2,4\},\{1,3,4\},\{2,3,4\}\}$ and $\ws\TSL\TSC(\TSP)= \WSLC\text{*}(\TSP, \tau) = \{\varphi, \TSP ,\{1\}, \{2\}, \{3\},\{4\},\{1,2\},\{1,3\}, \{1,4\}, \{2,3\},~\{2,4\},~\{3,4\},~\{1, 2, 3\}~\{1,2,4\},\break \{1,3,4\},\{2,3,4\}\}$, and
$\WSLC\text{{*}{*}}(\TSP, \tau)= \{\varphi, \TSP ,\{1\}, \{2\}, \{3\},\{4\},\{1,2\},\{1,3\},\break \{1,4\}, \{2,3\}, \{2,4\},\{3,4\}, \{1, 2,3\},\{1,3,4\},\{2,3,4\}\}$ and locally closed sets or $\LC = \{\varphi, \TSP, \{1\}, \{2\}, \{4\}, \{1,2\}, \{3, 4\}\{1,2,3\},\{1,3,4\},\{2,3,4\}\}$, Here $\{3\}$ is $\WSLC (\TSP, \tau)$, $\WSLC\text{*}(\TSP, \tau)$, $\WSLC\text{{*}{*}}(\TSP, \tau)$ but not LC-set.
\end{exm}

\begin{thm}\label{thm7.2.2}
If a subset $\clrD$ of $(\TSP, \tau)$ is $\WSLC*$-set then it is $\WSLC$-set but inverse is untrue.
\end{thm}

\begin{proof}
Take up $\clrD$ be a $\WSLC*$-set. Assume $\TSM$ as a ws-open set in $\TSP$ and $\TSN$ as closed set in $\TSP$. Seeing that $\clrD$ is $\ws\LC*$-set by definition, $\clrD = \TSM \cap \TSN$. Since each closed set is ws-closed. Henceforth $\clrD$ is $\WSLC$-set.

We use example \ref{exam7.2.3} to prove the inverse of theorem is untrue.
\end{proof}

\begin{exm}\label{exam7.2.3}
$\TSP = \{1, 2, 3, 4\}$ and $\tau = \{\varphi, \{1\},\{1, 2\}, \{1, 2, 3\}\}$ be a topology on $\TSP$. Then $\ws\TSC(\TSP) = \{\varphi, \TSP, \{2\}, \{3\}, \{4\},\{1, 4\} \{3,4\}, \{2,3\}, \{2,4\},\{1,2,4\},\{1,3,4\},\break \{2,3,4\}\} \ws\TSO(\TSP) = \{\varphi, \TSP, \{1\}, \{2\}, \{3\},\{1,2\}, \{1,3\},\{3,4\}, \{1,4\}, \{2,3\}, \{1,2,3\},\break \{1,2,4\},\{1,3,4\},\{2,3,4\}\}$ and $\ws\TSL\TSC(\TSP)= \{\varphi, \TSP ,\{1\}, \{2\}, \{3\},\{4\},\{1,2\},\{1,3\},\break \{1,4\},\{2,3\}, \{2,4\},\{3,4\}, \{1, 2, 3\}\{1,2,4\},\{1,3,4\},\{2,3,4\}\}$ And $\ws\TSL\TSC\text{*}(\TSP)=\break \{\varphi, \TSP ,\{1\},\{2\}, \{3\},\{4\},\{1,2\},\{1,3\}, \{1,4\}, \{2,3\},\{3,4\},~\{1, 2, 3\},\{1,2,4\},~\{1,3,4\},\break\{2,3,4\}\}$, Then the subset $\{2,4\} \in \WSLC(\TSP, \tau)$ but $\{2,4\}\not\in \WSLC\text{*}(\TSP, \tau)$.
\end{exm}

\begin{thm}\label{thm7.2.3}
Each $\WSLC${*}{*}$(\TSP, \tau)$ is $\WSLC (\TSP, \tau)$ but inverse is untrue.
\end{thm}

\begin{proof}
Proof of above theorem follows from definition \ref{defi7.2.3}.

We use example \ref{exam7.2.4} to prove the inverse of theorem is untrue.
\end{proof}

\begin{exm}\label{exam7.2.4}
$\TSP = \{1, 2, 3, 4\}$ and $\tau = \{\varphi, \TSP, \{1\},\{1, 2\}, \{1, 2, 3\}\}$ be a topology on $\TSP$. Then $\ws\TSC(\TSP) = \{\varphi, \TSP, \{2\}, \{3\}, \{4\},\{1, 4\} \{3,4\}, \{2,3\}, \{2,4\},\{1,2,4\},\{1,3,4\},\break\{2,3,4\}\} \ws\TSO(\TSP) = \{\varphi, \TSP, \{1\}, \{2\}, \{3\},\{1,2\}, \{1,3\},\{3,4\}, \{1,4\}, \{2,3\}, \{1,2,3\},\break \{1,2,4\},\{1,3,4\},\{2,3,4\}\}$ and $\ws\TSL\TSC(\TSP)= \{\varphi, \TSP ,\{1\}, \{2\}, \{3\},\{4\},\{1,2\},\{1,3\},\break \{1,4\}, \{2,3\}, \{2,4\},\{3,4\}, \{1, 2, 3\}\{1,2,4\},\{1,3,4\},\{2,3,4\}\}$ and $\WSLC\text{{*}{*}}(\TSP, \tau)= \{\varphi, \TSP,\{1\},\{2\}, \{3\},\{4\},\{1,2\},\{1,4\}, \{1,3\}, \{1,4\},\{3,4\}, \{1, 2, 3\},\{1,2,4\},\{1,3,4\},\break\{1,3,4\}\}$. Here $\{1,3\} \in \WSLC (\TSP, \tau)$, but $\{1, 3\}\not\in \WSLC\text{{*}{*}}(\TSP, \tau)$.
\end{exm}

\begin{thm}\label{thm7.2.4}
Each locally closed set is a wslc set but inverse is untrue.
\end{thm}

\begin{proof}
Proof is obtained from the facts that each open set is ws-open and each closed set is ws-closed.
\end{proof}

\begin{exm}\label{exam7.2.5}
$\TSP = \{1, 2, 3, 4\}$ and $\tau = \{\varphi, \{1\}, \{2\}, \{1, 2\}, \{1, 2,3\}\}$ be a topology on $\TSP$. Then $\ws\TSC(\TSP) = \{\varphi, \TSP,\{1\}, \{2\}, \{3\}, \{4\}, \{1, 3\}, \{1, 4\} \{3,4\},\{2,3\}, \{2,4\},\break \{1,2, 4\},\{1,3,4\},\{2,3,4\}\}, \ws\TSO(\TSP) = \{\varphi, \TSP, \{1\}, \{2\}, \{3\},\{1,2\}, \{1,3\}, \{1,4\},\{2,3\},\break \{2,4\}, \{1,2, 3\}, \{1,2,4\},\{1,3,4\},\{2,3,4\}\}$ and $\ws\TSL\TSC(\TSP)= \{\varphi, \TSP ,\{1\}, \{2\},\{3\},\{4\},\break \{1,2\},\{1,3\}, \{1,4\}, \{2,3\}, \{2,4\},\{3,4\}, \{1, 2, 3\}\{1,2,4\},\{1,3,4\},\{2,3,4\}\}$, and\break locally closed sets or $\LC = \{ \varphi, \TSP, \{1\}, \{2\}, \{4\}, \{1,2\}, \{3, 4\} \{1,2,3\},\{1,3,4\},\break \{2,3,4\}\}$. Here $\{3\}$ is $\WSLC (\TSP, \tau)$, but not LC-set.
\end{exm}

\begin{thm}\label{thm7.2.5}
Each $\alpha - \lc(g\xi\text{*-}\lc, \alpha \gp\text{-}\lc$, $*g\alpha$-$\lc$) set is a $\wslc$ set but the inverse is untrue.
\end{thm}

\begin{proof}
Proof is obtained from the fact that each $\alpha-(\TSg\xi\text{*}, \alpha \gp\text{-}, \text{*}\TSg\alpha)$ closed set is ws-closed and each $\alpha -(\TSg\xi\text{*}, \alpha\gp, \text{*}\TSg\alpha)$ open set is ws-open set.
\end{proof}

\begin{exm}\label{exam7.2.6}
$\TSP = \{1, 2, 3, 4\}$ and $\tau = \{\varphi, \{1\}, \{2\}, \{1, 2\}, \{1, 2,3\}\}$ be a topology on $\TSP$. Then $\ws\TSC(\TSP) = \{\varphi, \TSP, \{1\}, \{2\}, \{3\}, \{4\}, \{1, 3\}, \{1, 4\} \{3,4\},\{2,3\}, \{2,4\},\break \{1,2, 4\},\{1,3,4\},\{2,3,4\}\}$, $\ws\TSO(\TSP) = \{\varphi, \TSP, \{1\}, \{2\}, \{3\},\{1,2\}, \{1,3\}, \{1,4\},\break \{2,3\}, \{2,4\}, \{1,2, 3\}, \{1,2,4\},\{1,3,4\},\{2,3,4\}\}$ and $\ws\LC(\TSP)= \{\varphi, \TSP ,\{1\}, \{2\},\break \{3\},\{4\},\{1,2\},\{1,3\}, \{1,4\}, \{2,3\}, \{2,4\},\{3,4\},\{1, 2, 3\},\{1,2,4\},\{1,3,4\},\{2,3,4\}$\break and $\alpha$-$\lc(g\xi\text{*-}\lc, \alpha\TSg\TSp\text{-}\lc, \text{*}\TSg\alpha\text{-}\lc) = \{\varphi, \TSP, \{1\}, \{2\}, \{3\}, \{4\}, \{1,2\}, \{3, 4\} \{1,2,3\},\break \{1,2,4\},\{1,3,4\},\{2,3,4\}\}$. Here $\{2,3\}$ is $\ws\LC(\TSP, \tau)$, but not $\alpha\text{-}\lc(\TSg\xi\text{*-}\lc$, $\alpha\TSg\TSp\text{-}\lc$, *g$\alpha$-$\lc$) set.
\end{exm}

\begin{thm}\label{thm7.2.6}
Every $\TSg^{\text{\#}}$-$\lc$(resp, $\ddot{g}$-$\lc$) set is a wslc set but the converse is not true.

\end{thm}

\begin{proof}
The proof follows from the fact that every $g^{\text{\#}}$-closed(resp, $\ddot{g}$-closed) set is ws-closed and every $g^{\text{\#}}$-open(resp, $\ddot{g}$-open) set is ws-open set.
\end{proof}

\begin{exm}\label{exam7.2.7}
Let $\TSP = \{1, 2, 3, 4\}$ and $\tau = \{\varphi, \{1\}, \{1, 2\}, \{1, 2, 3\}\}$ be a topology on $\TSP$. Then $\ws\TSC(\TSP) = \{\varphi, \TSP, \{2\}, \{3\}, \{4\},\{1, 4\} \{3,4\}, \{2,3\}, \{2,4\},\{1,2,4\},\break \{1,3,4\},\{2,3,4\}\} \ws\TSO(\TSP) = \{\varphi, \TSP, \{1\}, \{2\}, \{3\},\{1,2\}, \{1,3\},\{3,4\}, \{1,4\}, \{2,3\},\break \{1,2,3\}, \{1,2,4\},\{1,3,4\},\{2,3,4\}\}$ and $\ws\TSL\TSC(\TSP)= \{\varphi, \TSP,\{1\}, \{2\}, \{3\},\{4\},\{1,2\},\break \{1,3\}, \{1,4\},\{2,3\}, \{2,4\},\{3,4\}, \{1, 2, 3\}\{1,2,4\},\{1,3,4\},\{2,3,4\}\}$ and $g^{\text{\#}}$-$\lc(\ddot{g}\text{-}\lc)\break = \{\varphi,\TSP,\{1\},\{4\},\{1,2\},\{3, 4\},\{1,2,3\},\{2,3,4\}\}$. Here $\{2\}$ is $\ws\LC(\TSP, \tau)$, but not $g^{\text{\#}}\text{-}\lc(\ddot{g}\text{-}\lc)$ set.
\end{exm}

\begin{thm}\label{thm7.2.7}
Every $\scl\text{-}\lc$ set is a $\wslc$ set but the converse is not true.
\end{thm}

\begin{proof}
The proof follows from the fact that every semi closed set is ws-closed and every semi - open set is ws-open set.
\end{proof}

\begin{exm}\label{exam7.2.8}
Let $\TSP = \{1, 2, 3, 4\}$ and $\tau = \{\varphi, \{1\}, \{1, 2\}, \{1, 2, 3\}\}$ be a topology on $\TSP$. Then $\ws\TSC(\TSP) = \{\varphi, \TSP, \{2\}, \{3\}, \{4\},\{1, 4\} \{3,4\}, \{2,3\}, \{2,4\},\{1,2, 4\},\break \{1,3,4\},\{2,3,4\}\}, \ws\TSO(\TSP) = \{\varphi, \TSP, \{1\}, \{2\}, \{3\},\{1,2\}, \{1,3\},\{3,4\}, \{1,4\},\{2,3\},\break \{1,2, 3\},\{1,2,4\},\{1,3,4\},\{2,3,4\}\}$ and $\ws\LC(\TSP)= \{\varphi, \TSP,\{1\}, \{2\}, \{3\},\{4\},\{1,2\},\break\{1,3\}, \{1,4\},\{2,3\}, \{2,4\},\{3,4\}, \{1, 2, 3\}\{1,2,4\},\{1,3,4\},\{2,3,4\}\}$ and $\scl\text{-}\lc=\break \{ \varphi, \TSP ,\{1\}, \{2\},\{3\},\{4\},\{1,2\},\{1,3\}, \{1,4\}, \{2,3\}, \{2,4\},\{3,4\},\{1,2,4\},~\{1,3,4\},\break \{2,3,4\}\}$. Here $\{1,2,3\}$ is $\ws\LC (\TSP, \tau)$, but not $\scl\text{-}\lc$ set.
\end{exm}

\begin{thm}\label{thm7.2.8}
\begin{enumerate}[(1)]
\item $\wslc$ set is $\gsprlc$ set but not conversely.
\item $\wslc$ set is $\gsplc$ set but not conversely.
\item $\wslc$ set is $\rgblc$ set but not conversely
\end{enumerate}
\end{thm}

\begin{proof}
The proof follows from the two definitions and the fact that every ws-closed (resp. ws-open) set is $\gspr$-closed (resp. $\gspr$-open).

Similrly we can prove the result (2) and (3)
\end{proof}

\begin{exm}\label{exam7.2.9}
Let $\TSP= \{1,2, 3\}$ and $\tau=\{\TSP, \phi, \{1\}, \{2,3\}\}$. Then $\{2\}$ is a gsprlc(respectively, gsplc, rgblc) set but not a ws-lc set in $(\TSX, \tau)$.
\end{exm}

\begin{rem}\label{rem7.2.2}
The following example shows that wlc, $\hat{g}$lc, $\alpha\hat{g}$lc, pgprlc, swglc, g$\alpha$lc, gw$\alpha$lc, $\alpha$gslc and rwlc are independent with wslc sets.
\end{rem}

\begin{exm}\label{exam7.2.10}
Let $\TSP= \{1, 2, 3, 4\}$ and $\tau=\{\TSP, \phi, \{1\}, \{2\}, \{1, 2\}, \{1, 2, 3\}\}$. Then $\{1, 2\}$ is a wslc set but not a w-lc (respectively, $\hat{g}$lc, $\alpha\hat{g}$lc, pgprlc, swglc, g$\alpha$lc ,gw$\alpha$lc, $\alpha$gslc and rwlc) set in $(\TSP, \tau)$.
\end{exm}

\begin{exm}\label{exam7.2.11}
Let $\TSP= \{1,2,3, 4\}$ and $\tau=\{\TSP, \phi, \{1, 2\}, \{3,4\}\}$. Then $\{1\}$ is w-lc (respectively, $\hat{g}$lc, $\alpha\hat{g}$lc, pgprlc, swglc, g$\alpha$lc, gw$\alpha$lc, $\alpha$gslc and rwlc) set but not a wslc set in $(\TSP, \tau)$.
\end{exm}

\begin{rem}\label{rem7.2.3}
From the above discussions and known facts we have the following implications:
\begin{figure}[H]
\centering
\includegraphics[scale=.75]{1.eps}
\end{figure}
\end{rem}

\begin{thm}\label{thm7.2.9}
In a space $(\TSP, \tau)$, if each semi closed set is closed then semi $\LC(\TSP) \subset \ws\LC(\TSP)$
\end{thm}

\begin{proof}
In general, $\LC(\TSP) \subseteq \ws\LC(\TSP)$ and $\LC(\TSP) \subseteq \semi\LC(\TSP)$. Since each semi closed set is closed then semi-$\LC(\TSP) \subset \LC(\TSP)$. Then semi-$\LC(\TSP) = \LC(\TSP)$. Therefore semi-$\LC(\TSP) \subset \semi\LC(\TSP)$.
\end{proof}

\begin{thm}\label{thm7.2.10}
The given results are one and the same.
\begin{enumerate}[(1)]
\item $\clrD \in\ws\LC\text{*}(\TSP, \tau)$
\item $\clrD = \TSM \cap \cl(\clrD)$ for some ws open set $\TSM$
\item $\cl(\clrD)\text{-}\clrD$ is ws closed
\item $\clrD \cup \cl(\clrD)^{\text{c}}$ is ws open, where $\clrD$ is subset of a topological space $(\TSP, \tau)$
\end{enumerate}
\end{thm}

\newpage

\begin{proof}
(1) $\Rightarrow$ (2)

Take up $\clrD \in \ws\LC\text{*}\TSP$. Then $\clrD = \TSM \cap \TSF$, where $\TSM$ and $\TSF$ are ws open set and a closed set in $\TSP$. Seeing that $\clrD \subseteq \TSF$, $\cl(\clrD) \subseteq \cl(\TSF) = \TSF$. Now $\TSM \cap \cl(\clrD) \subseteq (\TSM \cap \TSF) = \TSD$. That is $\TSM \cap \cl(\clrD) \subseteq \TSD$. Also, $\clrD \subseteq \TSM$ and $\clrD \subseteq\cl(\clrD)$ then $\clrD \subseteq (\TSM \cap \cl(\clrD))$. Therefore $\TSD= \TSM \cap \cl(\clrD)$ for some ws open set $\TSM$.

(2) $\Rightarrow$ (1)

Here $\TSM$ is ws open and $\TSF$ is ws closed set. Therefore $\clrD = (\TSM \cap \cl(\clrD)) \in \ws\LC * (\TSP)$. (iii) $\Rightarrow$ (iv) Take up $\cl(\clrD) - \clrD$ is ws-closed. Now $\TSP - (\cl(\clrD) - \clrD) = \TSP \cap (\cl(\clrD) - \TSD)^{\text{c}} = \clrD \cup (\TSP - \cl(\clrD)) = \clrD \cup (\cl(\clrD))^{\text{c}}$ is a ws-open set, That is $\clrD \cup \cl(\clrD)^{\text{c}}$ is ws open set.

(4) $\Rightarrow$ (3)

Take up $\clrD \cup \cl(\clrD)^{\text{c}}$ is ws open. Now $\TSP - (\clrD \cup (\cl(\clrD))^{\text{c}}) = \cl(\clrD) \cap (\TSP - \TSD) = \cl(\clrD) - \clrD$ is ws-closed.

(4) $\Rightarrow$ (2)

Take up $\clrD \cup \cl(\clrD)^{\text{c}}$ is ws open. Now $\TSM\cap \cl(\clrD) = (\clrD \cup (\cl(\clrD))^{\text{c}}) \cup \cl(\clrD) = (\clrD \cap \cl(\clrD)) \cup ((\cl(\clrD))^{\text{c}} \cap \cl(\clrD)) = \clrD \cup \varphi = \clrD$

(2) $\Rightarrow$ (4)

Take up $\clrD = \TSM \cap \cl(\clrD)$ for some ws open set $\TSM$. Now $\clrD \cup (\cl(\clrD)^{\text{c}}) = (\TSM \cap \cl(\clrD)) \cup (\cl(\clrD)^{\text{c}}) = \TSM \cap \cl(\clrD) \cup (\cl(\clrD))^{\text{c}} = \TSM \cap \TSP = \TSM$ is a ws-open set.
\end{proof}

\begin{thm}\label{thm7.2.11}
If $\clrD \in \ws\LC(\TSP, \tau)$ then $\clrD = \TSM\cap \ws\cl(\clrD)$ for some open set $\TSM$, where $\clrD$ is the subset of a topological space $(\TSP, \tau)$
\end{thm}

\begin{proof}
Take up $\clrD \in \ws\LC(\TSP)$. Then $\clrD = \TSM \cap \TSF$ where $\TSM$ and $\TSF$ are ws closed and ws open sets in $(\TSP,\tau)$ respectively. Since $\clrD \subseteq \TSF$, $\ws\cl(\clrD) \subseteq \ws\cl(\TSF)=\TSF$. Now $\TSM \cap \ws\cl(\clrD) \subseteq (\TSM \cap \TSF) = \TSD$. Also, since $\clrD \subseteq \TSM$ and $\clrD \subseteq \ws\cl(\clrD)$ then $\clrD \subseteq (\TSM \cap \ws\cl(\clrD))$. Therefore $\clrD = \TSM \cap \ws\cl(\clrD)$.
\end{proof}

\begin{thm}\label{thm7.2.12}
If $\TSA \in \ws\LC{*}{*}(\TSP, \tau)$ then $\clrD = \TSM \cap \ws\cl(\clrD)$, where $\TSM$ is an open set and $\clrD$ is subset of a topological space $(\TSP, \tau)$
\end{thm}

\begin{proof}
Take up $\clrD \in \ws\LC{*}{*}(\TSP, \tau)$. Then $\clrD = \TSM \cap \TSF$, where $\TSM$ is an open set and $\TSF$ is a ws-closed set. Seeing that $\clrD \subseteq \TSF$ and $\ws\cl(\clrD) \subseteq \ws\cl(\TSF) = \TSF$ thus $\TSM \cap \ws\cl(\clrD) \subseteq (\TSM \cap \TSF) = \TSD$. That is $\TSM \cap \ws\cl(\clrD) \subseteq \clrD$. Also $\clrD \subseteq \TSM$ and $\clrD \subseteq \ws\cl(\clrD)$ then $\clrD \subseteq \TSM \cap \ws\cl(\clrD)$. Thus, $\clrD = \TSM \cap \ws\cl(\clrD)$ for some open set $\TSM$.
\end{proof}

\section{WSLC-Continuous Maps}\label{sec7.3}

\begin{dfn}\label{defi7.3.1}
The map $\TSh: \TSP \to \TSQ$ is called a $\WSLC$-continuous (respectively $\WSLC$*-continuous and $\WSLC${*}{*}-continuous) if the inverse image of each open set in $\TSQ$ is $\ws\lc$-set(respectively $\ws\lc$*-set and $\ws\lc${*}{*}-set) in $(\TSP, \tau)$.
\end{dfn}

\begin{exm}\label{exam7.3.1}
Take up $\TSP = \TSQ = \{1, 2, 3\}$, $\tau = \{\phi, \{1\}, \{2\}, \{1, 2\}\}$ and $\sigma = \{\phi, \TSQ, \{1\}, \{2,3\}\}$. Then $\WSLC(\TSP) = \TSP (\TSP)$ and $\TSh: \TSP \to \TSQ$ is defined as identity map is a $\WSLC$-continuous.
\end{exm}

\begin{thm}\label{thm7.3.1}
Take up $\TSh: (\TSP) \to \TSQ$ be a function. Then If $\TSH$ is locally continuous then it is $\WSLC$-continuous but inverse is untrue
\end{thm}

\begin{proof}
Proof is obvious from the well known definitions and evidence
\end{proof}

\begin{exm}\label{exam7.3.2}
Take $\TSP = \{1, 2, 3, 4\}$ and $\tau = \{\phi, \TSP \{1\}, \{2\}, \{1, 2\}, \{1, 2, 3\}\}$ be a topology on $\TSP$. $\sigma = \{\phi, \TSY, \{1\},\{ 2, 3\}\}$ Then $\ws\TSC(\TSP) = \{\phi, \TSP,\{1\}, \{2\}, \{3\}, \{4\}, \{1,3\},\break \{1, 4\} \{3,4\},\{2,3\}, \{2,4\}, \{1,2, 3\},\{1,3,4\},\{2,3,4\}\}$ $\ws\TSO(\TSP) = \{\phi, \TSP, \{1\}, \{2\}, \{3\},\break \{1,2\}, \{3,4\}, \{2,4\},\{2,3\}, \{2,4\}, \{1,2,3\}, \{1,2,4\},\{1,3,4\},\{2,3,4\}\}$ and $\ws\LC(\TSP)= \{\phi,\TSP ,\{1\}, \{2\},\{3\},\{4\},\{1,2\},\{1,3\}, \{1,4\}, \{2,3\},~ \{2,4\},\{3,4\}, \{1,2,3\},\{1,2,4\},\break \{1,3,4\},\{1,3,4\}\}$, $\LC(\TSP)=\{ \phi \TSP,\{1\}, \{2\}, \{4\}, \{1, 2\}, \{3,4\}, \{1,2,3\},\{1,3,4\},\break \{2,3,4\}\}$ and Let $\TSh: (\TSP, \tau )\to(\TSY, \sigma)$ be a function defined by identity map is $\ws\lc$-continuous but not a $\LC$ continuous function as the open set $\TSF=\{2,3\}$ in $\TSQ, h^{-1} (\TSF)=\{2,3\}$ is not a locally continuos closed set in $\TSP$.
\end{exm}

\begin{thm}\label{thm7.3.2}
Each $\alpha$-$\lc$ continuous ($\TSg\xi$*-lc continuous, $\alpha\gp$-$\lc$ continuous, $*\TSg\alpha$-$\lc$ continuous) is a $\ws\lc$-continuous but the inverse is untrue.
\end{thm}

\begin{proof}
Proof is obvious from the well known definitions and evidence.
\end{proof}

\begin{exm}\label{exam7.3.3}
Take $\TSP = \{1, 2, 3, 4\}$ and $\TSQ = \{1, 2, 3\}$, $\tau = \{\phi, \{1\}, \{2\}, \{1, 2\},\break \{1, 2, 3\}\}$ be a topology on $\TSP$. $\sigma = \{\phi, \TSQ, \{1\},\{ 2, 3\}\}$ be a topology on $\TSQ, \ws\TSC(\TSP) = \{\phi, \TSP,\{1\}, \{2\}, \{3\}, \{4\}, \{1,3\}, \{1, 4\} \{3,4\}, \{2,3\}, \{2,4\}, \{1,2, 3\},\{1,3,4\},\{2,3,4\}\}$,\break $\ws\TSO(\TSP) = \{\phi, \TSP, \{1\}, \{2\},\{3\},\{1,2\}, \{3,4\}, \{1,4\}, \{2,3\}, \{2,4\}, \{1,2,3\}, \{1,2,4\},\break \{1,3,4\},\{2,3,4\}\}$ and $\ws\LC(\TSP)= \{\phi, \TSP,\{1\}, \{2\}, \{3\},\{4\},\{1,2\},\{1,3\}, \{1,4\},\break \{2,3\}, \{2,4\},\{3,4\}, \{1,2,3\}\{1,2,4\},\{1,3,4\},\{2,3,4\}\}$, and $\alpha - \lc(\TSg\xi* -\lc, \alpha \gp-\lc, *\TSg\alpha-\lc) = \{ \phi, \TSP, \{1\}, \{2\}, \{3\}, \{4\}, \{1,2\}, \{3, 4\} \{1,2,3\},\{1,2,4\},\{1,3,4\},\break \{2,3,4\}\}$, Here $\{2,3\}$ is $\ws\LC (\TSP, \tau)$, but not $\alpha - \lc(\TSg\xi* -\lc, \alpha\gp-\lc, *\TSg\alpha-\lc)$ set. Let $\TSh: (\TSP, \tau )\to(\TSY, \sigma)$ be a function defined by identity map is $\ws\lc$-continuous but not a $\alpha$-$\lc$ continuous ($\TSg\xi*$-$\lc$ continuous, $\alpha\gp$-$\lc$ continuous, *$\TSg\alpha$-$\lc$ continuous)) function as the open set $\TSF=\{2,3\}$ in $\TSQ, h^{-1} (\TSF)=\{2,3\}$ is not a $\alpha - \lc(\TSg\xi * -\lc, \alpha\gp-\lc, *\TSg\alpha-\lc)^{\text{c}}$ losed set in $\TSP$.
\end{exm}

\begin{thm}\label{thm7.3.3}
Each semi-$\lc$-continuous ($g^{\text{\#}}$-$\lc$ continuos, $\ddot{g}$-$\lc$ continuos) set is a $\ws\lc$ continuos, but the inverse is untrue.
\end{thm}

\begin{proof}
Proof is obvious from the well known definitions and evidence.
\end{proof}

\begin{exm}\label{exam7.3.4}
Take $\TSP = \{1,2,3,4\}$ and $\TSQ= \{1,2,3\}$ and $\tau = \{\phi, \{1\},\{2\}, \{1,2\},\break \{1,2,3\}\}$ be a topology on $\TSP$. $\sigma = \{\phi, \TSQ, \{1\},\{ 2, 3\}\}$ be a topology on $\TSP$. $\ws\TSC(\TSP) = \{\phi, \TSP,\{1\}, \{2\}, \{3\}, \{4\},\{1,3\}, \{1, 4\} \{3,4\}, \{2,3\}, \{2,4\}, \{1,2, 3\},\{1,3,4\},\{2,3,4\}\}$\break $\ws\TSO(\TSP) = \{\phi, \TSP, \{1\}, \{2\}, \{3\},\{1,2\},\{3,4\}, \{1,4\}, \{2,3\}, \{2,4\}, \{1,2,3\}, \{1,2,4\},\break \{1,3,4\},\{2,3,4\}\}$ and $\ws\LC(\TSP)= \{\phi, \TSP ,\{1\}, \{2\},\{3\},\{4\},\{1,2\},\{1,3\}, \{1,4\},\break \{2,3\}, \{2,4\},\{3,4\}, \{1,2,3\}\{1,2,4\},\{1,3,4\},\{2,3,4\}\}$, and $\scl$-$\lc(g^{\text{\#}}-\lc, \ddot{g}-\lc) = \{\phi, \TSP,\{1\},\{4\},\{1,2\}\{3, 4\},\{1,2,3\},\{2,3,4\}\}$. Let $\TSh: (\TSP, \tau )\to(\TSY, \sigma)$ be a function defined by identity map is $\ws\lc$-continuous but not a semilc-continuous($g^{\text{\#}} -\lc$ continuos, $\ddot{g}$-$\lc$ continuos) function as the open set $\TSF=\{2,3\}$ in $\TSQ, h^{-1} (\TSF)=\{2,3\}$ is not a $\scl$-$\lc(g^{\text{\#}} -\lc, \ddot{g} -\lc)^{\text{c}}$ losed set in $\TSP$.
\end{exm}

\begin{thm}\label{thm7.3.4}
Every $\ws$-$\lc$ continuous is $\gspr$ (respectively $\gsp$-$\lc$ continuos, $\rgb$-$\lc$ continuos) continuos, but the converse is not true.
\end{thm}

\begin{proof}
The proof follows from the known definitions and facts.
\end{proof}

\begin{exm}\label{exam7.3.5}
$\TSX = \{1, 2, 3\}$ and $\TSY= \{1, 2, 3\}$ and $\tau = \{\TSX, \phi, \{1\}, \{2, 3\}\}$ be a topology on $\TSX. \sigma = \{\phi, \TSY, \{1\}, \{1, 2\}, \{1, 2, 3\}\}$, Then $\ws\TSC(\TSX) = \{\phi, \TSX, \{1\}, \{2, 3\}\}$. and $\ws\LC(\TSX)= \{\phi, \TSX \{1\}, \{2, 3\}\}$, and $\gspr - \lc(\gsp-\lc, \rgb-\lc) = \{\TSP(\TSX)\}$ Let $\TSh: (\TSX, \tau)\to(\TSY, \sigma)$ be a function defined by $\TSh(1)= 1$, $\TSh(2)=2$, $\TSh(3)=3$, is $\gspr$-$\lc$(respectively, $\gsp$-$\lc$, $\rgb$-$\lc$)-continuous but not a $\ws$-$\lc$ continuos function as the open set $\TSH=\{1,2,3\}$ in $\TSQ$, $h^{-1}(\TSH)=\{1,2,3\}$ is not a $\ws$-$\lc$ set in $\TSX$.
\end{exm}

\begin{thm}\label{thm7.3.5}
Each $\WSLC${*}{*} continuous is $\WSLC$-continuous.but inverse is untrue.
\end{thm}

\begin{proof}
Proof is obvious from the well known definitions and evidence.
\end{proof}

\begin{exm}\label{exam7.3.6}
$\TSP = \{1,2,3,4\}$ and $\TSY= \{1,2,3\}$ and $\tau = \{\phi, \{1\},\{2\}, \{1,2\},\break \{1,2,3\}\}$ be a topology on $\TSP$. $\sigma = \{\phi, \TSY, \{1\},\{ 2, 3\}\}$, $\ws\TSC(\TSP) = \{\phi, \TSP,\{1\}, \{2\}, \{3\},\break \{4\}, \{1,3\}, \{1, 4\} \{3,4\}, \{2,3\}, \{2,4\}, \{1,2, 3\},\{1,3,4\},\{2,3,4\}\}$, $\ws\TSO(\TSP) = \{\phi, \TSP,\break \{1\}, \{2\}, \{3\},\{1,2\}, \{3,4\}, \{1,4\}, \{2,3\}, \{2,4\}, \{1,2,3\}, \{1,2,4\},\{1,3,4\},\{2,3,4\}\}$\break and $\ws\LC(\TSP)= \{\phi, \TSP ,\{1\}, \{2\}, \{3\},\{4\},\{1,2\},\{1,3\}, \{1,4\}, \{2,3\}, \{2,4\},\{3,4\},\break \{1,2,3\}\{1,2,4\},\{1,3,4\},\{1,3,4\}\}$ and $\ws\LC{*}{*}\TSP)= \{\phi, \TSP,\{1\}, \{2\}, \{3\},\{4\},\{1,2\},\break \{1,3\}, \{1,4\}, \{2,3\}, \{2,4\},\{3,4\}, \{1, 2, 3\},\{1,3,4\},\{2,3,4\}\}$, Let $\TSh: (\TSP, \tau )\to(\TSQ, \sigma)$ be a function defined by $\TSh(1)= 2$, $\TSh(2)=2$, $\TSh(3)=4$, is $\ws\lc$-continuous but not a $\ws\LC${*}{*}-continuous function as the open set $\TSH=\{2,3\}$ in $\TSQ$, $h^{-1} (\TSH)=\{1,2,3\}$ is not a $\ws\LC${*}{*} closed set in $\TSP$.
\end{exm}

\begin{rem}\label{rem7.3.1}
From the above discussions and known facts we have the following implications:
\begin{figure}[H]
\centering
\includegraphics[scale=.78]{2.eps}
\end{figure}
\end{rem}

\begin{thm}\label{thm7.3.6}
Let $\TSh: \TSP \to \TSQ$ is a function defined on a door space $(\TSP,\tau)$ is a $\WSLC$-continuous (respectively $\WSLC*$ continuous and $\WSLC${*}{*}-continuous).
\end{thm}

\begin{proof}
Take up $\TSh$ be a function, where $\TSQ$ be any space and $(\TSP, \tau)$ is a door space. Let $\TSE \in \sigma$. Then $\TSh^{-1} (\TSE)$ is either open or closed. In both the cases, $\TSh^{-1} (\TSE) \in \WSLC(\TSP, \tau)$ (respectively $\WSLC$*$(\TSP)$, $\WSLC${*}{*}$(\TSP)$). Therefore $\TSh$ is $\WSLC$-continuous (respectively $\WSLC$* continuous and $\WSLC${*}{*}-continuous).
\end{proof}

\begin{thm}\label{thm7.3.7}
Take $\TSh : \TSP \to \TSQ$ be $\WSLC$-continuous and $\TSE$ is an open set in $\TSQ$. Then $\TSh : (\TSP, \tau) \to (\TSE, \tau \TSE)$ is $\WSLC$-continuous.
\end{thm}

\begin{proof}
Take up $\TSU$ as an open set in $\TSE$ and hence open in $(\TSQ, \sigma)$. Then $\TSh^{-1} (\TSV)$ is $\ws\lc$-set in $(\TSP, \tau)$, by hypothesis. Therefore $\TSh : (\TSP, \tau) \to (\TSE, \tau \TSE)$ is $\WSLC$-continuous.
\end{proof}

\begin{thm}\label{thm7.3.8}
Take $\TSh: \TSP \to \TSQ$ is $\WSLC$-continuous (respectively $\WSLC$*-{\break}continuous and $\WSLC${*}{*}-continuous) and $\TSg: \TSQ \to \TSR$ is continuous then $\goh : \TSP \to \TSR$ is $\WSLC$-continuous (respectively $\WSLC*$-continuous and $\WSLC${*}{*}-continuous).
\end{thm}

\begin{proof}
Take up $\TSN$ as any open set in $(\TSR, \eta)$. Seeing that $\TSg$ is continuous map, $\TSg^{-1} (\TSN)$ is open in $(\TSQ, \sigma)$. Also $\TSh$ is $\WSLC$-continuous, $\TSh^{-1} (\TSg^{-1} (\TSN)) = (\goh)^{-1} (\TSN)$ is a $\ws\lc$ set in $(\TSP, \tau)$. Then $\goh$ is $\WSLC$-continuous.
\end{proof}

\section{WSLC irresolute maps and their properties}\label{sec7.4}

\begin{dfn}\label{defi7.4.1}
A map $\TSh: \TSP \to \TSQ$ is called $\ws\lc$-irresolute (respectively $\WSLC$*-irresolute, $\WSLC${*}{*}-irresolute) if inverse image of each $\ws$ locally closed set in $\TSQ$ is ws locally closed set (recpectively $\ws\lc${*} locally closed, $\ws\lc${*}{*} locally closed) set in $(\TSP, \tau)$.
\end{dfn}

\begin{exm}\label{exam7.4.1}
Take up $\TSP = \TSQ = \{1,2, 3\}$, $\tau = \{\phi, \{1\}, \{2\}, \{1,2\}\}$ and $\sigma = \{\phi, \TSQ, \{1\}, \{2, 3\}\}$. Then $\WSLC(\TSP) = \TSP (\TSP)$ and $\TSh: \TSP \to \TSQ$ is defined as identity map is a $\WSLC$-irresolute map.
\end{exm}

\begin{thm}\label{thm7.4.1}
If $\TSh: (\TSP) \to \TSQ$ is $\ws$-irresolute then $\TSh$ is $\WSLC$-irresolute but inverse is untrue.
\end{thm}

\begin{proof}
Take up $\TSh : (\TSP) \to \TSQ$ is $\ws$-irresolute. Let $\TSK \in \WSLC(\TSP)$. Then $\TSK = \TSL \cap \TSM$, where $\TSL$ is $\ws$-open set and $\TSM$ is $\ws$-closed set in $(\TSQ, \sigma)$. Now, $\TSh^{-1} (\TSK) = \TSh^{-1} (\TSL \cap \TSM) = \TSh^{-1} (\TSL) \cap \TSh^{-1} (\TSK)$. Seeing that $\TSh$ is $\ws$-irresolute, $\TSh^{-1}(\TSL)$ is $\ws$-open and $\TSh^{-1}(\TSM)$ is $\ws$ closed set in $(\TSP, \tau)$, therefore $\TSh^{-1}(\TSM)\in \ws\lc(\TSP)$, Hence the theorem.
\end{proof}

\begin{exm}\label{exam7.4.2}
Take up $\TSP = \TSQ = \{1, 2, 3, 4\}$, and $\tau= \{\phi, \TSP, \{1\}, \{1, 2\}, \{1, 2, 3\}\}$. $\sigma = \{\phi, \TSQ, \{1\}, \{2\}, \{1, 2\}, \{1, 2, 3\}\}$, $\ws\TSC(\TSP) = \{\phi, \TSP, \{2\}, \{3\}, \{4\}, \{1, 4\}, \{2,3\},\break \{2,4\},\{3,4\}, \{1,2, 4\},\{1,3,4\},\{2,3,4\}\}$, $\ws\TSC(\TSQ) = \{\phi, \TSQ,\{1\}, \{2\}, \{3\}, \{4\}, \{1, 3\},\break \{1, 4\} \{2,3\}, \{2,4\}, \{3,4\}\{1,2, 4\},\{1,3,4\},\{2,3,4\}\}$. Let $\TSh : \TSP \to \TSQ$ be an identity map but not a $\ws$-irresolute map, as the inverse image of the $\ws$-closed set $\{1\}$ in $\TSQ$ is $\{1\}$ in $(\TSP, \tau)$, which is not a $\ws$-closed set in $(\TSP, \tau)$.
\end{exm}

\begin{thm}\label{thm7.4.2}
If a map $\TSh: (\TSP) \to \TSQ$ is $\WSLC$-irresolute then it is $\ws\lc$-continuous.
\end{thm}

\begin{proof}
Result is obvious from the well known definitions and evidence that each open set is a $\ws\lc$-set
\end{proof}

\begin{thm}\label{thm7.4.3}
Any function $\TSh : (\TSP) \to \TSQ$ defined on a door space is $\WSLC$-irresolute.
\end{thm}

\begin{proof}
The proof is obvious.
\end{proof}

\begin{thm}\label{thm7.4.4}
Consider two functions $\TSh: \TSP \to \TSQ$ and $\TSg: \TSQ\to (\TSR)$. Then
\begin{itemize}
\item[(i)] $\goh : (\TSP, \tau) \to (\TSR, \eta)$ is $\WSLC$-irresolute if both $\TSh$ and $\TSg$ are $\WSLC$-irresolute maps.
\item[(ii)] $\goh: (\TSP, \tau) \to (\TSR, \eta)$ is $\WSLC$-continuous if $\TSh$ is $\WSLC$-continuous and $\TSg$ is $\WSLC$-irresolute.
\end{itemize}
\end{thm}

\begin{proof}
Take up $\TSL$ be any $\ws\lc$ set in $(\TSR, \eta)$. By hypothesis, $\TSg^{-1} (\TSL)$ is a $\ws\lc$ set in $(\TSQ, \sigma)$. Also $\TSh$ is $\WSLC$-irresolute, $\TSh^{-1} (\TSg^{-1} (\TSL)) = (\goh)^{-1} (\TSL)$ is a $\ws\lc$ set in $(\TSP, \tau)$. Therefore goh is $\WSLC$-irresolute.

This theorem also hold for $\WSLC$*-irresolute map and $\WSLC${*}{*}-irresolute map. (ii) Take up $\TSL$ be any open set in $(\TSR, \eta)$. By hypothesis, $\TSg^{-1} (\TSL)$ is a $\ws\lc$ set in $(\TSQ, \sigma)$. Also, since $\TSh$ is $\WSLC$-irresolute, $\TSh^{-1} (\TSg^{-1} (\TSL)) = (\goh)^{-1} (\TSL)$ is a $\ws\lc$ set in $(\TSP, \tau)$. Therefore $\goh$ is $\WSLC$-continuous.
\end{proof}

\begin{dfn}\label{defi7.4.2}
Function $\TSh: \TSP \to \TSQ$ is said to be sub-$\WSLC$*-continuous if there is a basis $\beta$ for $\TSQ\ni: \TSh^{-1} (\TSV) \in \WSLC\text{*}(\TSP)$ for each $\TSV \in \beta$.
\end{dfn}

\begin{thm}\label{thm7.4.5}
If $\TSh : (\TSP, \tau) \to \TSQ$ is $\WSLC$*-continuous then it is sub-$\WSLC$*-continuous.
\end{thm}

\begin{proof}
Take up $\TSh: \TSP \to \TSQ$ be $\WSLC*$-continuous and $\beta$ is a bais for $\TSQ$ and $\TSu \in \beta$. Since $\TSh$ is $\WSLC$*-continuous, $\TSh^{-1}(\TSu) \in \WSLC \text{*} (\TSP, \tau)$. Therefore $\TSh$ is sub-$\WSLC$*-continuous.
\end{proof}

\begin{thm}\label{thm7.4.6}
If $\TSh: \TSP \to \TSQ$ is sub-$\LC$-continuous then it is sub-$\WSLC$*-continuous.
\end{thm}

\begin{proof}
Result is obvious from the well known definitions and evidence that $\LC(\TSP,\tau) \subseteq \WSLC(\TSP,\tau)$.
\end{proof}

\begin{thm}\label{thm7.4.7}
Take up $\clrD$ and $\TSE$ be two subsets of a topological space $(\TSP, \tau)$. Then these statements are verified.
\begin{enumerate}[(i)]
\item If $\clrD \in \WSLC \text{*} (\TSP)$ and $\TSE \in \WSLC \text{*} (\TSP)$ then $(\TSD \cap \TSE) \in \WSLC \text{*} (\TSP)$.
\item If $\clrD \in \WSLC\text{{*}{*}}(\TSP)$ and $\TSE$ is open then $(\TSD \cap \TSE) \in \WSLC\text{{*}{*}}(\TSP)$.
\item If $\clrD \in \WSLC(\TSP)$ and $\TSE$ is $\ws$-open then $(\TSD \cap \TSE) \in \WSLC(\TSP)$
\item If $\clrD \in \WSLC\text{*}(\TSP)$ and $\TSE$ is $\ws$-open then $(\TSD \cap \TSE) \in \WSLC\text{*}(\TSP)$.
\item If $\clrD \in \WSLC\text{*}(\TSP)$ and $\TSE$ is closed then $(\TSD \cap \TSE) \in \WSLC\text{*}(\TSP)$.
\end{enumerate}
\end{thm}

\begin{proof}
\begin{itemize}
\item[(i)] Take up $\clrD \in \WSLC\text{*}(\TSP)$ and $\TSE \in \WSLC\text{*}(\TSP)$. and also we know that $\exists$ two $\ws$-open sets $\TSK$ and $\TSL \ni: \clrD = \TSK \cap \cl(\clrD)$ and $\TSE = \TSL \cap \cl(\TSE)$. Now $\TSD \cap \TSE = (\TSK \cap \TSL) \cap (\cl(\clrD) \cap \cl(\TSE))$. Here $\TSK \cap \TSL$ is $\ws$-open and $(\cl(\clrD) \cap \cl(\TSE))$ is closed. Therefore $(\TSD \cap \TSE) \in \WSLC\text{*}(\TSP)$.

\item[(ii)] Take up $\clrD \in \WSLC\text{{*}{*}}(\TSP)$ and $\TSE$ is open. Then $\clrD = \TSK \cap \TSG$, where $\TSK$ is an open set and $\TSG$ is a $\ws$-closed. Consider $\TSD \cap \TSE = (\TSK \cap \TSE) \cap \TSG$, where $\TSK \cap \TSE$ is open and $\TSG$ is $\ws$-closed. Therefore $(\TSD \cap \TSE) \in \WSLC\text{{*}{*}}(\TSP)$.

\item[(iii)] Take up $\clrD \in \WSLC(\TSP)$ and $\TSE$ is $\ws$-open. Then $\clrD = \TSK \cap \TSG$, where $\TSP$ is $\ws$-open and $\TSG$ is $\ws$-closed. Consider $\TSD \cap \TSE = (\TSK \cap \TSE) \cap \TSG$ where $\TSK \cap \TSE$ is $\ws$-open and $\TSG$ is $\ws$-closed. Therefore $(\TSD \cap \TSE) \in \WSLC(\TSP)$.
\end{itemize}
Similarly we can prove (iv) and (v).
\end{proof}
